% ---------------------------------------------------------------------
% Chapter: cross-model comparison (Llama-3.2-3B, Qwen2.5-7B, Qwen2.5-14B)
% Numbers are macros from ../results/tables/numbers.tex. No literals.
% ---------------------------------------------------------------------

\chapter{The Scaling Curve: Three Models Compared}
\label{ch:comparison}

Chapter~\ref{ch:qwen} established each larger model on its own terms. This
chapter puts all three on one axis --- Llama-3.2-3B, Qwen2.5-7B and
Qwen2.5-14B, over an identical corpus, an identical question set and identical
policy code --- and asks what changes with scale, what does not, and which of
those changes were expected. Table~\ref{tab:scaling-results} carries the full
numbers; the analysis here reads them.

\section{What Is and Is Not Comparable}
\label{sec:cmp-validity}

Holding the corpus, questions and code fixed controls almost everything, but
not the one variable the central claim depends on. ``Selective retrieval beats
always-retrieve'' is a statement about accuracy \emph{at a given retrieval
budget}, and the three gated policies do not share a budget: they retrieve on
\retPfiveMcq{}, \retQwenPfiveMcq{} and \retQwenFourteenPfiveMcq{} of
multiple-choice queries respectively. The budgets fall as the model grows, for
the reason set out in Section~\ref{sec:qwen-budget}: the hallucination probe is
threshold-free on multiple-choice and cannot be recalibrated, and a larger
model self-agrees more often, so the ensemble's realised rate drifts downward
with scale.

\begin{quote}
  Within each model, the policies are compared at a controlled budget and the
  ordering is meaningful. \emph{Across} models the gated rows sit at different
  budgets, so a raw difference between two P5 rows confounds model scale with
  retrieval rate. The comparisons that survive are those that fix the budget by
  construction --- the closed-book policies, which all retrieve 0\% --- and the
  \emph{within-model} orderings, which are read as directions rather than
  cross-model magnitudes.
\end{quote}

\section{Closed-Book Capability: A Curve That Flattens}
\label{sec:cmp-mcq}

Closed-book accuracy is the cleanest cross-model comparison, since no policy
retrieves and no budget confound applies. It rises then plateaus:
\accPthreeMcq{} at 3B, \accQwenPthreeMcq{} at 7B, \accQwenFourteenPthreeMcq{}
at 14B. The two steps are not equal. The 3B-to-7B step adds
\scaleGainSmall{}pp; the 7B-to-14B step, across a model twice the size, adds
only \scaleGainLarge{}pp. Doubling the parameters a second time buys almost
nothing on this benchmark.

For the 3B-to-7B step, both models answered the same \mcnemarN{} questions, so
a paired test is available: \mcnemarBoth{} correct for both, \mcnemarNeither{}
for neither, \mcnemarOnlyQwen{} for Qwen alone and \mcnemarOnlyLlama{} for
Llama alone, giving $\chi^2 = \mcnemarChi{}$ on the discordant pairs
($p < 0.05$). The \mcnemarOnlyLlama{} questions the smaller model gets right
and the larger gets wrong are a reminder that scale redistributes competence
rather than simply adding it; a single aggregate accuracy is not a complete
description of a model's medical knowledge. From 7B to 14B the aggregate barely
moves, which is the more consequential fact for this dissertation and is taken
up in Section~\ref{sec:cmp-bars}.

\section{Retrieval's Value Inverts, Then Stays Inverted}
\label{sec:cmp-open}

The open-ended arm is complete at $n=200$ for every policy and model, once the
gated runs used thresholds refitted for open-ended drafts
(Section~\ref{sec:qwen-saturation}). The closed-book scores rise and plateau in
the same shape as multiple-choice: \fOnePthreeOpen{}, \fOneQwenPthreeOpen{},
\fOneQwenFourteenPthreeOpen{}.

The finding is in the \emph{ordering} of the three policies, and how it changes
with scale:

\begin{quote}
  \textbf{3B:} P1 (\fOnePoneOpen{}) $>$ P5 (\fOnePfiveOpen{}) $>$ P3
  (\fOnePthreeOpen{}) --- always-retrieve is best, retrieval helps.\\
  \textbf{7B:} P3 (\fOneQwenPthreeOpen{}) $>$ P5 (\fOneQwenPfiveOpen{}) $>$ P1
  (\fOneQwenPoneOpen{}) --- closed-book is best, retrieval hurts.\\
  \textbf{14B:} P3 (\fOneQwenFourteenPthreeOpen{}) $>$ P5
  (\fOneQwenFourteenPfiveOpen{}) $>$ P1 (\fOneQwenFourteenPoneOpen{}) ---
  the same inverted order as 7B.
\end{quote}

Two things matter here. First, the ordering \emph{inverts} between the smallest
model and the larger two: on the 3B, open-ended questions sit outside the
model's parametric knowledge, so retrieved context is a net help and more
retrieval scores higher; by 7B the model already knows more than the passages
reliably add, so the same corpus that helped now hurts. Second, the inversion
does not drift back --- it holds identically at 14B. The reversal is therefore
not a quirk of the 7B but a stable property of the two larger models, which is
the strongest form the observation can take without a fourth scale. Open-ended
questions are simply where parametric knowledge runs out last, so the
retrieval-hurts regime that multiple-choice showed at every scale reaches
open-ended only once the model is large enough.

The usual cautions bound the gated rows: the open-ended budgets are not matched
across models (\retPfiveOpen{}, \retQwenPfiveOpen{}, \retQwenFourteenPoneOpen{}
are not the operating points to compare), so the P5 figures are read as ordinal
positions within each model's own P1--P3 range; and paragraph-level token-$F_1$
is surface-lexical, so only orderings are interpreted, never the size of a gap.

\section{The Mechanism: Distraction Scales with Capability}
\label{sec:cmp-mechanism}

Why does retrieval hurt more as the model improves? The paired
closed-book-versus-always-retrieve analysis, run at each scale, answers it
directly. Counting the multiple-choice questions where always-retrieving
\emph{breaks} a closed-book-correct answer against those it \emph{fixes}:

\begin{quote}
  3B: \breakLlama{} broken / \fixLlama{} fixed (\harmRatioLlama{}$\times$)
  \quad 7B: \breakQwenSeven{} / \fixQwenSeven{}
  (\harmRatioQwenSeven{}$\times$) \quad
  14B: \breakQwenFourteen{} / \fixQwenFourteen{}
  (\harmRatioQwenFourteen{}$\times$).
\end{quote}

The fixes never vanish --- the corpus keeps supplying answers the model lacked
at every scale --- but the breakages grow, and the 14B breaks the most correct
answers of any model. This is the signature of displacement, not of a
deficient corpus: a corpus that simply lacked information would damage a
stronger model \emph{less}, because the stronger model can fall back on what it
knows, whereas here the strongest model is harmed the most precisely because it
has the most correct knowledge to be displaced. Section~\ref{sec:disc-synthesis}
develops this into the central reframing --- that a gate's value is protection,
not efficiency --- and rules out the corpus-absence objection in full. For the scaling story it is enough to note that the phenomenon the gate
exists to manage does not weaken with scale; it strengthens.

\section{Selective Beats Always-Retrieve --- and at 14B, Significantly}
\label{sec:cmp-selective}

The central claim is that gating beats always-retrieving. Within each model,
where the comparison is legitimate, P5 leads P1 on multiple-choice at all three
scales: by \diffPfivePone{}pp on the 3B, \diffQwenPfivePone{}pp on the 7B and
\diffFourteenPfivePone{}pp on the 14B. On the two smaller models the 95\%
intervals (\diffPfivePoneCI{} and \diffQwenPfivePoneCI{}) cross zero: each is
individually \emph{consistent with} an improvement rather than a demonstration
of one. At 14B the interval is \diffFourteenPfivePoneCI{} --- it excludes zero,
so for the first time the advantage of selective over always-retrieve is
statistically significant on a single model.

That significance is a direct consequence of the mechanism in
Section~\ref{sec:cmp-mechanism}: because distraction is worst on the 14B,
always-retrieve is punished hardest there, opening the widest and cleanest gap
for the gate to exploit. The claim is thus supported two ways at once --- by
three same-direction estimates whose agreement across scales is itself evidence,
and by one estimate that clears significance on its own at the scale where the
effect is strongest.

\section{Scale Versus the Safety Bars}
\label{sec:cmp-bars}

The most consequential cross-model fact is where the curve sits relative to the
clinical bars (\barLow{}/\barMedium{}/\barHigh{}). The 3B cleared none. Both
Qwen models clear the \barLow{} low-risk bar on their best policy
(\accQwenPthreeMcq{} and \accQwenFourteenPthreeMcq{}), so the safety envelope
that was empty at 3B is populated at low risk once capability is sufficient. But
\emph{neither larger model reaches the \barMedium{} medium-risk bar}, and the
curve is flattening as it approaches: the \scaleGainLarge{}pp gained from 7B to
14B is a twelfth of the \scaleGainSmall{}pp gained from 3B to 7B. Extrapolating
this curve to the 80\% bar would require a scale increase that the measured
diminishing returns give no reason to expect to be enough.

This is the result that most sharpens the dissertation's argument. Scaling is a
real lever --- it moved a system from clearing no bar to clearing the low one
--- but it is a saturating lever, and it saturates below the threshold that
medium-risk clinical use would demand. That is precisely the regime in which
\emph{method} matters: if raw capability cannot be scaled to safety within a
practical range, then the discipline of retrieving only when it helps, and not
otherwise, is not an efficiency nicety but part of how the remaining distance is
to be closed.

\section{Was This the Expected Result?}
\label{sec:cmp-expected}

Partly, and the unexpected parts are the more useful. Three predictions held.
Closed-book capability rose with scale and the paired test supported it at
$p<0.05$. Calibration non-transfer, predicted from each smaller model's signal
distribution, was confirmed prospectively at every new scale. And selective
retrieval beat always-retrieve in direction at all three scales, becoming
significant on the 14B.

Three things were not predicted. The closed-book curve \emph{plateaus} sharply
after 7B, so ``just use a bigger model'' is a weaker remedy than it first
appears. The medium-risk bar stays \emph{out of reach} even at 14B, which
relocates the burden from capability onto method. And the open-ended
retrieval-helps result from the 3B did not merely fail to replicate but
\emph{inverted}, then held inverted at 14B --- a sharper and more stable finding
than a clean replication would have been. Together they say that the value of
retrieval is not a fixed property of the corpus or the method but a function of
the gap between what the model knows and what the passages add: a gap that
scaling closes on the tasks the model can already do, and that the gate exists
to detect on the tasks it cannot. That is the thesis the project was built to
test, and it took the third model to turn it from a two-point line into a curve
with a shape.

\section{Cost}
\label{sec:cmp-cost}

Latency is not comparable across all three models, because the 3B ran on CPU
and both Qwen models on a T4 GPU; that comparison would measure the accelerator,
not the model. The within-GPU comparison is clean, however, now that the gated
runs were re-profiled with the heavy trackers disabled. On the T4, closed-book
answers a multiple-choice question in \latQwenPthreeMcq{} (7B) and
\latQwenFourteenPthreeMcq{} (14B); the gated policy takes \latQwenPfiveMcq{} and
\latQwenFourteenPfiveMcq{} respectively. The gate's cost is the price of its
extra draft generations and scales gently with model size, as expected --- the
14B's gated query is slower than the 7B's but remains within an interactive
budget. Cost therefore does not disturb the accuracy conclusions, and the
deployability argument (Section~\ref{sec:disc-deployment}) rests on a measured
gate cost rather than an estimated one.
